\subsection{Objekt M18}
\label{subsec:ObjektM18}

\small

\paragraph{1. Objektsteckbrief}

\par Bei dem Objekt M18 handelt es sich um ein Mehrfamilienwohnhaus im Marterlweg 18, 92334 Pollanten. Es handelt sich um ein im Jahr 1987 errichtetes Gebäude. Das ursprüngliche \gls{efh} wurde in den Jahren 2023 bis 2025 umfassend grundsaniert und in ein \gls{mfh} mit 3 Wohneinheiten umgebaut. Das Objekt verfügt über 3 Wohneinheiten mit einer vermietbaren Fläche von insgesamt $284{}m^2$ Wohnen verteilt auf 3 WE. Die Vermietungsquote liegt konstant bei 100\%. Das Objekt befindet sich im Eigentum der \gls{stz} es besteht jedoch ein Nießbrauchrecht des Komplementärs und der Komanditistin Angelika Seitz.


\begin{table} [htbp!]
\begin{tabular*}{\textwidth}{@{\extracolsep{\fill}}>{\raggedright\arraybackslash}p{4.5cm}>{\raggedright\arraybackslash}p{10.3cm}}
\toprule
\textbf{Merkmal} & \textbf{Angabe} \\
\midrule
Adresse & Marterlweg 18, 92334 Pollanten \\
Objektart & \gls{mfh} \\
Baujahr / Modernisierung & 1988 / 2023 - 2025 \\
Vermietbare Fläche & $284 m^2$ Wohnen \\
Einheiten & 3 Wohneinheiten \\
Vermietungsstand & -  per \today \\
Eigentümer & \gls{stz} \\
\bottomrule
\end{tabular*}
\caption{Objektsteckbrief M18}
\label{tab:objektsteckbriefM18}
\end{table}

\paragraph{2. Ertragslage und Stabilität (p.a.)}
\begin{table} [htbp!]
\begin{tabular*}{\textwidth}{@{\extracolsep{\fill}}>{\raggedright\arraybackslash}p{8.0cm}>{\raggedleft\arraybackslash}p{2.2cm}>{\raggedleft\arraybackslash}p{2.2cm}>{\raggedleft\arraybackslash}p{2.2cm}}
\toprule
\textbf{Kennzahl} & \textbf{2023} & \textbf{2024} & \textbf{2025} \\
\midrule
Nettokaltmiete 	& 0 		&  0 	&  0 \\
Nicht umlagefähige Bewirtschaftungskosten & - & - & - \\
NOI (Net Operating Income) & - & - & - \\
Leerstandsquote & 0\% & 0\% & 0\% \\
\bottomrule
\end{tabular*}
\caption{Ertragslage M18}
\label{tab:ertragslageM18}
\end{table}


\paragraph{3. Bankrelevante Kennzahlen (stabilisierter Zustand)}
\begin{table} [htbp!]
\begin{tabular*}{\textwidth}{@{\extracolsep{\fill}}>{\raggedright\arraybackslash}p{6.1cm}>{\raggedleft\arraybackslash}p{2.8cm}>{\raggedright\arraybackslash}p{6.3cm}}
\toprule
\textbf{Kennzahl} & \textbf{Wert} & \textbf{Einordnung} \\
\midrule
Marktwert / interner Ansatz & [EUR] & Bewertungsstichtag: [Datum] \\
Beleihungswert (konservativ) & [EUR] & Sicherheitsabschlag: [\%] \\
Wunschdarlehen auf Objektbasis & [EUR] & Anteil am Gesamtobligo: [\%] \\
LTV (Darlehen / Marktwert) & [\%] & Bankziel: [z.B. \textless= 70\%] \\
Debt Yield (NOI / Darlehen) & [\%] & Bankziel: [z.B. \textgreater= 8\%] \\
DSCR (CFADS / Schuldendienst) & [x,xx] & Bankziel: [z.B. \textgreater= 1,20] \\
\bottomrule
\end{tabular*}
\caption{Bankrelevante Kennzahlen T5}
\label{tab:bankkennzahlenT5}
\end{table}

\paragraph{4. Risiko- und Sicherheitenprofil}
\begin{itemize}
	\item \textbf{Mietausfall-/Leerstandsrisiko:} Da mit Nießbrauchrecht der Komanditistin und des Komplementärs kein Leerstandsrisiko, da keine Vermietung angedacht instand
	\item \textbf{Instandhaltungsrisiko:} Da grundlegend Saniert und in Nießbrauch kein Risiko für die \gls{kg}
	\item \textbf{Marktrisiko:} Lage in Pollanten mit stabilem Immobilienmarkt 
	\item \textbf{Rechtliche Sicherheiten:} 
\end{itemize}

\paragraph{5. Kurzfazit und Ausblick}
Das Bestandsobjekt M18 zeigt einen sehr hervorragenden stabilen Zustand. Der erfolgreiche Umbau des Objekts unter Federführung der Eigentümerin der \gls{kg} unter dem Dach der verbundenen \gls{m18} kann als Blaupause für künftige Sanierungsprojekte gewertet werden.  Unter konservativen Annahmen ergibt sich ein belastbarer Kreditrahmen auf das Objekt M18 von \textbf{[EUR]}, der durch werthaltige Sicherheiten und ein transparentes Risikomanagement unterlegt ist.

\vspace{0.5em}
\textit{Hinweis: Detailrechnungen (Tilgungsplan, Sensitivitäten, Mieterliste) sind im Anhang bzw. in den Tabellenblättern hinterlegt.}

\normalsize
